\begin{solucion}{Costo adicional del minorista}
Definimos:
\[
\begin{array}{c|l}
\toprule
\text{Simbolo} & \text{Significado} \\
\midrule
q\in\mathbb{R}_+ & \text{cantidad vendida.}\\
p(q)=a-bq & \text{precio final inducido por la demanda inversa.}\\
w\in\mathbb{R}_+ & \text{precio mayorista.}\\
c\in\mathbb{R}_+ & \text{costo unitario del proveedor.}\\
k\in\mathbb{R}_+ & \text{costo operativo unitario del minorista.}\\
\bottomrule
\end{array}
\]
El costo \(k\) no lo recibe el proveedor; es un costo real adicional que reduce el margen unitario del minorista.

Con costo operativo \(k\), la utilidad del minorista es:
\[
\Pi_r(q)=q(p-w-k)
=q(a-bq-w-k).
\]
Derivando:
\[
\frac{d\Pi_r}{dq}=a-w-k-2bq=0,
\]
por lo que la respuesta del minorista es:
\[
q(w)=\frac{a-w-k}{2b}.
\]

El proveedor maximiza:
\[
\Pi_s(w)=(w-c)q(w)
=(w-c)\frac{a-w-k}{2b}.
\]
Derivando:
\[
\frac{d\Pi_s}{dw}=0
\quad\Rightarrow\quad
w_D=\frac{a+c-k}{2}.
\]
Entonces:
\[
q_D=\frac{a-c-k}{4b},
\qquad
p_D=\frac{3a+c+k}{4}.
\]
Las utilidades son:
\[
\Pi_s^D=\frac{(a-c-k)^2}{8b},
\qquad
\Pi_r^D=\frac{(a-c-k)^2}{16b},
\]
\[
\Pi_{SC}^D=\frac{3(a-c-k)^2}{16b}.
\]

En la cadena centralizada, el costo marginal total es \(c+k\):
\[
\Pi_{SC}(q)=q(p-c-k)=q(a-bq-c-k).
\]
Por lo tanto:
\[
q_C=\frac{a-c-k}{2b},
\qquad
p_C=\frac{a+c+k}{2},
\]
\[
\Pi_{SC}^C=\frac{(a-c-k)^2}{4b}.
\]

Aplicando \(a=100\), \(b=1\), \(c=1\), \(k=5\):
\[
w_D=\frac{100+1-5}{2}=48,
\]
\[
q_D=\frac{100-1-5}{4}=23.5,
\qquad
p_D=\frac{3(100)+1+5}{4}=76.5.
\]
Utilidades descentralizadas:
\[
\Pi_s^D=\frac{94^2}{8}=1104.5,
\qquad
\Pi_r^D=\frac{94^2}{16}=552.25,
\]
\[
\Pi_{SC}^D=1656.75.
\]
Centralizado:
\[
q_C=\frac{94}{2}=47,
\qquad
p_C=\frac{100+1+5}{2}=53,
\]
\[
\Pi_{SC}^C=\frac{94^2}{4}=2209.
\]

El costo del minorista reduce la cantidad eficiente y las utilidades, pero la conclusion no cambia: la cadena descentralizada sigue vendiendo menos que la centralizada.
\end{solucion}
